% Oral Advanced Problem Solving - groupe G
% Compilation : pdflatex diapo.tex (depuis docs/oral/), ou import du dossier docs/oral dans Overleaf.
% Les figures sont dans figures/ et se régénèrent avec : python -m bench.plots
\documentclass[aspectratio=169,11pt]{beamer}

\usepackage[utf8]{inputenc}
\usepackage[T1]{fontenc}
\usepackage[french]{babel}
\usepackage{graphicx}
\usepackage{booktabs}

\usetheme{default}
\usecolortheme{dove}
\setbeamertemplate{navigation symbols}{}
\setbeamertemplate{footline}[frame number]
\setbeamerfont{frametitle}{series=\bfseries}
\graphicspath{{figures/}}

% \fig{nom}{hauteur} : insère figures/nom.png, ou un cadre vide si l'image n'existe pas encore
\newcommand{\fig}[2]{%
  \begin{center}
  \IfFileExists{figures/#1.png}%
    {\includegraphics[height=#2\textheight,width=\textwidth,keepaspectratio]{#1}}%
    {\fbox{\parbox[c][#2\textheight][c]{0.8\textwidth}{\centering figure à générer : #1.png}}}
  \end{center}}

% \afaire{texte} : rappel visible dans le PDF, à supprimer avant l'oral
\newcommand{\afaire}[1]{\textcolor{red}{[#1]}}

\title{Streaming Videos -- Hash Code 2017}
\subtitle{Analyse expérimentale de notre solution}
\author{Groupe G \\ \afaire{noms}}
\date{12 octobre 2026}

\begin{document}

% ---------------------------------------------------------------- 1
\begin{frame}
  \titlepage
\end{frame}

% ================================================================
\section{Algorithmique}

% ---------------------------------------------------------------- 2
\begin{frame}{Le problème}
  \begin{columns}
    \column{0.55\textwidth}
    \begin{itemize}
      \item Des vidéos, des caches de capacité limitée, des endpoints
      \item Chaque requête est servie par le cache le plus rapide qui contient la vidéo, sinon par le datacenter
      \item Score : latence moyenne économisée par requête
    \end{itemize}
    \column{0.45\textwidth}
    \afaire{schéma de l'énoncé}
  \end{columns}
\end{frame}

% ---------------------------------------------------------------- 3
\begin{frame}{Nos algorithmes}
  \small
  \begin{tabular}{@{}lll@{}}
    \toprule
    Algo & Idée & Particularité \\
    \midrule
    Glouton 1 & par requête, les plus grosses d'abord & cache le plus rapide qui a la place \\
    Glouton 2 & par vidéo, les plus populaires d'abord & cache vu par le plus de demandeurs \\
    Glouton 3 & par cache, endpoints par gain décroissant & une vidéo par endpoint \\
    \midrule
    Knapsack indépendant & un sac à dos par cache & doublons entre caches voisins \\
    Knapsack avec mise à jour & caches du plus au moins prometteur & valeurs revues après chaque cache \\
    \midrule
    Descente & meilleur voisin, 5 mouvements & part du meilleur glouton \\
    Tabou aléatoire & voisinage tiré au hasard & 100 itérations, liste de 7 \\
    Tabou trié & vidéos les plus lourdes en requêtes & déterministe \\
    \bottomrule
  \end{tabular}

  \medskip
  Sac à dos : deux DP exactes (sur le poids ou sur la valeur), repli sur un glouton au-delà de 10\,000.
\end{frame}

% ================================================================
\section{Évaluation expérimentale}

% ---------------------------------------------------------------- 4
\begin{frame}{Protocole expérimental}
  \begin{columns}[T]
    \column{0.5\textwidth}
    \textbf{Réglages}
    \begin{itemize}
      \item \afaire{CPU, RAM, version de Python : bench/data/machine.txt}
      \item 33 instances : 4 Google, 29 de la promo
      \item 8 algos, mêmes paramètres partout
      \item 10 graines pour le tabou aléatoire
    \end{itemize}
    \column{0.5\textwidth}
    \textbf{Limites}
    \begin{itemize}
      \item 300 s par exécution
      \item à la limite : on évalue l'état courant
    \end{itemize}
    \textbf{Métriques}
    \begin{itemize}
      \item écart au meilleur score connu (\%)
      \item temps d'exécution
      \item gain sur la solution de départ
      \item branche du sac à dos par cache
    \end{itemize}
  \end{columns}
\end{frame}

% ---------------------------------------------------------------- 5
\begin{frame}{Vue d'ensemble}
  \fig{vue_ensemble}{0.82}
\end{frame}

% ---------------------------------------------------------------- 6
\begin{frame}{Gloutons : lequel, et où ?}
  \fig{gloutons}{0.72}
  \afaire{une phrase de constat}
\end{frame}

% ---------------------------------------------------------------- 7
\begin{frame}{Gloutons : pourquoi ?}
  \begin{columns}
    \column{0.6\textwidth}
    \fig{capacite}{0.7}
    \column{0.4\textwidth}
    \begin{itemize}
      \item \afaire{cause 1}
      \item \afaire{cause 2}
      \item \afaire{cause 3}
    \end{itemize}
  \end{columns}
\end{frame}

% ---------------------------------------------------------------- 8
\begin{frame}{Recherche locale : le constat}
  \fig{recherche_locale}{0.72}
  \afaire{une phrase de constat}
\end{frame}

% ---------------------------------------------------------------- 9
\begin{frame}{Recherche locale : améliorable ?}
  \fig{voisinage}{0.62}
  \begin{itemize}
    \item Voisinage par défaut : 100 mouvements examinés par itération
    \item \afaire{effet d'un voisinage plus grand, coût en temps}
  \end{itemize}
\end{frame}

% ---------------------------------------------------------------- 10
\begin{frame}{Knapsack : quelle résolution pour quels caches ?}
  \fig{knapsack_branches}{0.78}
\end{frame}

% ---------------------------------------------------------------- 11
\begin{frame}{Knapsack : l'effet de la mise à jour}
  \begin{columns}
    \column{0.6\textwidth}
    \fig{knapsack_maj}{0.7}
    \column{0.4\textwidth}
    \begin{itemize}
      \item \afaire{gain moyen de la mise à jour}
      \item \afaire{lien avec le recouvrement des caches}
    \end{itemize}
  \end{columns}
\end{frame}

% ---------------------------------------------------------------- 12
\begin{frame}{SWOT}
  \small
  \begin{columns}[T]
    \column{0.5\textwidth}
    \begin{block}{Forces}
      \begin{itemize}
        \item \afaire{...}
      \end{itemize}
    \end{block}
    \begin{block}{Opportunités}
      \begin{itemize}
        \item \afaire{...}
      \end{itemize}
    \end{block}
    \column{0.5\textwidth}
    \begin{block}{Faiblesses}
      \begin{itemize}
        \item \afaire{...}
      \end{itemize}
    \end{block}
    \begin{block}{Menaces}
      \begin{itemize}
        \item \afaire{...}
      \end{itemize}
    \end{block}
  \end{columns}
\end{frame}

% ================================================================
\appendix

\begin{frame}{Secours : compromis temps / qualité}
  \fig{temps_qualite}{0.8}
\end{frame}

\begin{frame}{Secours : complexité des deux DP}
  \begin{itemize}
    \item DP sur le poids : $O(n \cdot W)$, $n$ vidéos utiles, $W$ capacité du cache
    \item DP sur la valeur : $O(n \cdot \sum v_i)$
    \item On prend la plus petite table ; si les deux dépassent 10\,000 : glouton par densité
    \item Poids et valeurs divisés par leur pgcd
  \end{itemize}
\end{frame}

\begin{frame}{Secours : corrections identifiées}
  \afaire{liste des corrections retenues}
\end{frame}

\end{document}
